\section{Stilbene Case Study: Long Molecular Sheet Evolution}
\label{sec:results-stilbene}

The stilbene dataset provides a larger molecular trajectory for testing whether
Reeb-space sheet tracking remains useful beyond the MVK transition case. We use
the same top-20 sheet selection and range-space shape-IoU correspondence score.
The strongest results are interval-level range events and long-lived sheet
families.

\subsection{Range-Space Event Intervals}

\begin{table}[t]
  \centering
  \caption{Important stilbene intervals identified by range-space Reeb-sheet
  tracking with threshold \(\theta=0.5\).}
  \label{tab:stilbene-events}
  \begin{tabular}{lrrrl}
    \toprule
    Interval & Score & Mean best & Weak source/target & Observation \\
    \midrule
    11915--11920 & 54.49 & 0.276 & 20/20 & strongest range reconfiguration \\
    11920--11925 & 53.44 & 0.328 & 20/20 & continuation of late event \\
    9380--9400 & 35.94 & 0.478 & 13/12 & earlier secondary event \\
    9660--9680 & 28.99 & 0.550 & 9/9 & secondary local reconfiguration \\
    11910--11915 & 28.13 & 0.494 & 9/9 & onset of late event \\
    2420--2440 & 27.65 & 0.518 & 9/9 & earlier secondary event \\
    \bottomrule
  \end{tabular}
\end{table}

After correcting duplicate timestep indices in the generated stilbene viewer
data, the dominant range-space event occurs at \(11915\)--\(11920\), followed by
\(11920\)--\(11925\). Unlike MVK, we do not currently tie these intervals to an
external application-level transition label; instead, they demonstrate that the
tracking pipeline identifies sheet reconfiguration intervals in a longer
molecular sequence.

\subsection{Long-Lived Feature Families}

The event intervals coexist with long-lived sheet tracks. The longest
threshold-\(0.5\) tracks include sheet \(86\rightarrow6270\) over labels
\(2100\)--\(5340\), with length 163 and mean continuation score 0.860; sheet
\(38\rightarrow19239\) over labels \(9400\)--\(11255\), with length 162 and mean
score 0.844; and sheet \(19\rightarrow53844\) over labels \(5540\)--\(8700\),
with length 159 and mean score 0.916. These tracks show that prominent
range-space sheets can remain coherent over long portions of the stilbene
trajectory even when short intervals exhibit strong reconfiguration.

\subsection{Stride And Threshold Sensitivity}

The \(11915\)--\(11920\) event is the top event for thresholds
\(\theta=0.3\) through \(\theta=0.7\). The sampling-stride check also preserves
the same late event: the top intervals at strides 1, 2, 3, and 4 are
\(11915\)--\(11920\), \(11910\)--\(11920\), \(11905\)--\(11920\), and
\(11900\)--\(11920\), respectively. Thus the late stilbene event is stable
under both threshold and timestep-stride changes.

\subsection{Metric Interpretation}

For stilbene, the combined score is identical to shape IoU. Shape IoU should
therefore remain the primary range metric. The best-target agreement fractions
against the combined score are 1.000 for shape IoU, 0.882 for bounding-box IoU,
0.849 for centroid similarity, and 0.568 for area ratio. Domain overlap agrees
with the shape-based target only 0.022 of the time. This supports using domain
overlap as a separate spatial-support view, not as a substitute for
range-space sheet tracking.

\subsection{Domain Events}

The domain event ranking provides a second view of the stilbene trajectory: it
identifies intervals where the same vertices stop supporting the same
Reeb-sheet behavior. The strongest domain events are \(8440\)--\(8460\),
\(8400\)--\(8420\), \(6940\)--\(6960\), \(8420\)--\(8440\), and
\(6920\)--\(6940\), each with weak domain continuations for all top sheets.
These intervals do not duplicate the range-event ranking; they indicate
separate domain-support reorganization events that should be inspected in the
domain view.

Domain/range target disagreement is now also useful for targeted visual checks:
the strongest corrected examples, such as \(5260\)--\(5280\),
\(8940\)--\(8960\), and \(9060\)--\(9080\), compare all 20 source sheets and
show low agreement between range-selected and domain-selected targets. We use
the domain event ranking as the main domain-space evidence and disagreement as
a way to choose concrete source/target examples.

\begin{figure*}[t]
  \centering
  \fbox{\parbox{0.92\linewidth}{
  Placeholder: stilbene range-event overview. Suggested panels: one around
  9380--9680 highlighting 9380--9400 and 9660--9680, and one around
  11900--11935 highlighting 11915--11920 and 11920--11925. Show the event graph
  and corresponding Sankey columns.}}
  \caption{Stilbene range-space event intervals identify a robust late sheet
  reconfiguration, with secondary earlier/local events.}
  \label{fig:stilbene-range-events}
\end{figure*}

\begin{figure}[t]
  \centering
  \fbox{\parbox{0.88\linewidth}{
  Placeholder: stilbene long-lived tracks. Highlight tracks
  \(86\rightarrow6270\), \(38\rightarrow19239\), and \(19\rightarrow53844\), and
  include clicked sheet/fiber detail views for one event interval.}}
  \caption{Long-lived stilbene Reeb-sheet tracks coexist with strong range
  event intervals.}
  \label{fig:stilbene-long-lived-tracks}
\end{figure}
